% ---------------- кодировки и язык ----------------
\usepackage[T2A]{fontenc}
\usepackage[utf8]{inputenc}
\usepackage[english,russian]{babel}

% ---------------- поля, интервалы, отступы ----------------
\usepackage[left=30mm,right=15mm,top=20mm,bottom=20mm]{geometry}
\usepackage{setspace}
\onehalfspacing
\usepackage{indentfirst}
\setlength{\parindent}{1.25cm}

% ---------------- математика, таблицы, списки ----------------
\usepackage{amsmath,amssymb}
\usepackage{array}
\usepackage{multirow}
\usepackage{enumitem}
\setlist{nosep}

% ---------------- рисунки и графики ----------------
\usepackage{graphicx}
\usepackage{float}
\usepackage{tikz}
\graphicspath{{img/}}

% ---------------- заголовки ----------------
\usepackage{titlesec}
\titleformat{\section}{\normalfont\Large}{\thesection}{1em}{}
\titleformat{\subsection}{\normalfont\large}{\thesubsection}{1em}{}
\titleformat{\subsubsection}{\normalfont\normalsize}{\thesubsubsection}{1em}{}

% ---------------- подписи ----------------
\usepackage{caption}
\captionsetup{labelsep=colon}
\addto\captionsrussian{%
	\renewcommand{\figurename}{Рис.}%
	\renewcommand{\tablename}{Таблица}%
	\renewcommand{\contentsname}{Содержание}%
	\renewcommand{\refname}{Список литературы}%
}

% ---------------- листинги ----------------
\usepackage{xcolor}
\usepackage{listings}
\renewcommand{\lstlistingname}{Листинг}
\lstset{
	language=Java,
	basicstyle=\footnotesize\ttfamily,
	keywordstyle=\color{blue},
	stringstyle=\color{red},
	commentstyle=\color{green!50!black},
	numbers=left,
	numberstyle=\tiny\color{gray},
	numbersep=6pt,
	backgroundcolor=\color{gray!5},
	frame=single,
	rulecolor=\color{gray!30},
	tabsize=4,
	breaklines=true,
	showstringspaces=false,
	captionpos=b,
	extendedchars=true,
	inputencoding=utf8,
	literate=
		{а}{{\cyra}}1 {А}{{\CYRA}}1 {б}{{\cyrb}}1 {Б}{{\CYRB}}1 {в}{{\cyrv}}1 {В}{{\CYRV}}1
		{г}{{\cyrg}}1 {Г}{{\CYRG}}1 {д}{{\cyrd}}1 {Д}{{\CYRD}}1 {е}{{\cyre}}1 {Е}{{\CYRE}}1
		{ё}{{\cyryo}}1 {Ё}{{\CYRYO}}1 {ж}{{\cyrzh}}1 {Ж}{{\CYRZH}}1 {з}{{\cyrz}}1 {З}{{\CYRZ}}1
		{и}{{\cyri}}1 {И}{{\CYRI}}1 {й}{{\cyrishrt}}1 {Й}{{\CYRISHRT}}1 {к}{{\cyrk}}1 {К}{{\CYRK}}1
		{л}{{\cyrl}}1 {Л}{{\CYRL}}1 {м}{{\cyrm}}1 {М}{{\CYRM}}1 {н}{{\cyrn}}1 {Н}{{\CYRN}}1
		{о}{{\cyro}}1 {О}{{\CYRO}}1 {п}{{\cyrp}}1 {П}{{\CYRP}}1 {р}{{\cyrr}}1 {Р}{{\CYRR}}1
		{с}{{\cyrs}}1 {С}{{\CYRS}}1 {т}{{\cyrt}}1 {Т}{{\CYRT}}1 {у}{{\cyru}}1 {У}{{\CYRU}}1
		{ф}{{\cyrf}}1 {Ф}{{\CYRF}}1 {х}{{\cyrh}}1 {Х}{{\CYRH}}1 {ц}{{\cyrc}}1 {Ц}{{\CYRC}}1
		{ч}{{\cyrch}}1 {Ч}{{\CYRCH}}1 {ш}{{\cyrsh}}1 {Ш}{{\CYRSH}}1 {щ}{{\cyrshch}}1 {Щ}{{\CYRSHCH}}1
		{ъ}{{\cyrhrdsn}}1 {Ъ}{{\CYRHRDSN}}1 {ы}{{\cyrery}}1 {Ы}{{\CYRERY}}1 {ь}{{\cyrsftsn}}1 {Ь}{{\CYRSFTSN}}1
		{э}{{\cyrerev}}1 {Э}{{\CYREREV}}1 {ю}{{\cyryu}}1 {Ю}{{\CYRYU}}1 {я}{{\cyrya}}1 {Я}{{\CYRYA}}1
		{—}{{\textemdash}}1 {→}{{$\to$}}1 {«}{{<<}}1 {»}{{>>}}1
}

% ---------------- ссылки ----------------
\usepackage{hyperref}
\hypersetup{
	colorlinks=true,
	linkcolor=blue,
	citecolor=green!60!black,
	urlcolor=blue,
	pdftitle={Отчет по лабораторной работе №1},
	pdfauthor={Храмов В.Д.}
}

% Поколение автомата в консольном виде (как печатает программа): файл figures/<имя>.txt.
% Ширина символа моноширинного шрифта 0.525em, клетка "# " = 1.05em, поэтому межстрочный
% интервал 1.05 размера шрифта делает клетки квадратными.
\lstdefinestyle{console}{
	language={}, numbers=none, frame=single, rulecolor=\color{gray!50},
	backgroundcolor=\color{gray!5}, framesep=2pt, xleftmargin=2pt, xrightmargin=2pt,
	breaklines=false, columns=fullflexible, keepspaces=true, showlines=true
}
\newcommand{\consolemedium}{\fontsize{9}{10.5}\selectfont\ttfamily}
\newcommand{\consolesmall}{\fontsize{7}{7.35}\selectfont\ttfamily}
\newcommand{\consoletiny}{\fontsize{5.2}{5.46}\selectfont\ttfamily}
\newcommand{\console}[2][\consolesmall]{\lstinputlisting[style=console,basicstyle=#1]{figures/#2.txt}}

% Клетка поля на рисунках (для схемы окрестности)
\newcommand{\nbcell}[3]{\draw[thick] (#1,#2) rectangle ++(1,1); \node at (#1+0.5,#2+0.5) {#3};}
